\documentclass[11pt]{article}
\usepackage[a4paper,margin=2.5cm]{geometry}
\usepackage[utf8]{inputenc}
\usepackage[T1]{fontenc}
\usepackage{lmodern}
\usepackage{amsmath,amssymb,amsthm}
\usepackage{booktabs}
\usepackage{pgfplots}
\usepgfplotslibrary{groupplots}
\usetikzlibrary{calc}
\usepackage[hidelinks]{hyperref}
\pgfplotsset{compat=1.18,
  every axis/.append style={width=0.9\linewidth, height=4.2cm, tick label style={font=\footnotesize},
    label style={font=\footnotesize}, legend style={font=\footnotesize, draw=none, fill=white, fill opacity=0.8, text opacity=1}},
  harvest/.style={fill=orange!25, draw=none},
  pathline/.style={blue!70!black, thin},
  thr/.style={red!70!black, dashed, thick}}
\newcommand{\datadir}{../sim/data}
\tikzset{thr/.style={red!70!black, dashed, thick}}


\newtheorem{theorem}{Theorem}[section]
\newtheorem{lemma}[theorem]{Lemma}
\newtheorem{proposition}[theorem]{Proposition}
\newtheorem{corollary}[theorem]{Corollary}
\theoremstyle{definition}
\newtheorem{assumption}[theorem]{Assumption}
\newtheorem{remark}[theorem]{Remark}

\newcommand{\R}{\mathbb{R}}
\newcommand{\dd}{\,\mathrm{d}}
\DeclareMathOperator{\sign}{sign}

\title{The optimal harvesting threshold\\ under a bounded harvesting rate is unique}
\author{Sebastian Senge\\Independent Researcher}
\date{}

\begin{document}
\maketitle

\begin{abstract}
A population in a random environment is harvested at a rate that cannot exceed a given capacity, and the goal is
the largest long-run average yield. For the stochastic logistic model it is known that an optimal strategy
harvests at full capacity above a threshold and not at all below it, and it was conjectured that the optimal
threshold is unique. We prove this.

For the stochastic logistic model the long-run yield is unimodal for all parameters under which the
unharvested population persists, so the optimal threshold exists and is unique. This includes the regime in which
harvesting at full capacity all the time does not let the population persist.

For one-dimensional diffusions whose total growth rate first increases and then decreases, and for a harvesting
capacity that may depend on the population size, a known down-crossing argument shows that the long-run yield of
a threshold strategy has at most one critical point; we state it without conditions on the boundary types.

The yield has a closed form in incomplete gamma functions, and bisection on an explicit sign function computes
the optimal threshold; it depends on two dimensionless numbers only.

For the logistic model, uniqueness, existence and unimodality, and with them the statement of the conjecture
about the yield function, are machine-checked in Lean~4 with Mathlib.
\end{abstract}

\section{Introduction}

Hunting and fishing remove animals from populations that fluctuate for reasons no manager controls: weather,
food supply, disease. Most harvest rules in practice combine two ingredients. One is a limit on effort or on the
quota, since the number of hunters, boats or licences is finite. The other is a population level below which
harvesting stops. In German game management, for instance, hoofed game other than wild boar may only be shot
under a cull plan that the authorities confirm or set, and the plan must be fulfilled \cite{bjagdg}; roe deer, red deer and
wild boar together account for close to two million animals per year \cite{djv}. Threshold rules of this kind
have a long history in the theory of fluctuating populations. In \cite{lande95} it was shown with
diffusion approximations that the strategies maximizing cumulative yield before extinction harvest at the
maximum possible rate above a threshold, and that a limited harvesting capacity puts that threshold below the
carrying capacity; threshold harvesting with uncertain population estimates was studied in \cite{lande97}.

A clean mathematical version of the question was posed in \cite{hening19}. The population follows the stochastic
logistic equation
\begin{equation}\label{eq:model}
  \dd X_t = X_t(\mu - \kappa X_t - h_t)\dd t + \sigma X_t \dd B_t ,\qquad 0 \le h_t \le M,
\end{equation}
where $\mu$ is the growth rate at low density, $\mu/\kappa$ the carrying capacity, $\sigma$ the strength of the
environmental noise, $h_t$ the per capita harvesting rate and $M$ its cap. The yield up to time $T$ is
$\int_0^T h_t X_t \dd t$, and the goal is the largest long-run average yield. When $\mu > \sigma^2/2$, which is
the condition for the unharvested population to persist, an optimal strategy is of threshold type
\cite[Theorem~2.1]{hening19}: it harvests at the full rate $M$ when $X_t$ exceeds a threshold $\eta$ and does not
harvest below it. The long-run yield $H(\eta)$ of the threshold $\eta$ is an explicit ratio of integrals
\cite[eq.~(4.4)]{hening19}. Whether the optimal threshold is unique was left open: it was conjectured in
\cite[Conjecture~4.1]{hening19}, with numerical support, that $H'(\eta) = 0$ has exactly one solution.

A related model with a harvesting quota that depends on the population size and an ergodic criterion in
expectation was studied in \cite{ekstrom26}, for diffusions with a natural boundary at $\infty$ and a cap that
does not decrease with the population size. There the first-order condition for the threshold has a unique root
\cite[Lemma~5.4]{ekstrom26}; uniqueness follows from a down-crossing argument, and existence uses an
integrability assumption that, for the logistic model with the cap $Mx$, amounts to $\mu - M > \sigma^2/2$. In
the logistic model $\infty$ is an entrance boundary \cite[Section~4]{hening19}, and the regime
$\mu - M \le \sigma^2/2$, in which harvesting at full capacity all the time does not let the population persist,
includes the example of Figure~1 of \cite{hening19}.
Related work treats a discounted criterion with a capped rate \cite{alvarez98}, unbounded harvesting
\cite{alvarezhening22}, harvesting combined with seeding \cite{heningtran20}, and discrete time
\cite{hening21}.

\paragraph{Contributions.}
\begin{itemize}
\item For the stochastic logistic model \eqref{eq:model}, Conjecture~4.1 of \cite{hening19} holds for all
$\mu > \sigma^2/2$, $\kappa > 0$, $M > 0$: the long-run yield has exactly one critical point $\eta^*$, it
increases strictly before $\eta^*$ and decreases strictly after it, and $\eta^*$ is the optimal threshold
(Theorem~\ref{thm:logistic}). The existence of the critical point when $\mu - M \le \sigma^2/2$ needs an
argument of its own (Lemma~\ref{lem:exist}).
\item The uniqueness argument of \cite{ekstrom26} holds for one-dimensional diffusions with drift $f$, noise
coefficient $v$ and a positive cap $k$ whenever $f$ first increases and then decreases, with no condition on the
types of the boundaries and no monotonicity of $k$, as long as the relevant integrals are finite and two boundary
terms vanish (Theorem~\ref{thm:general}).
\item The yield has a closed form in incomplete gamma functions. The optimal threshold equals
$\eta^* = \sigma^2\xi/(2\kappa)$, where $\xi$ depends only on $a = 2\mu/\sigma^2$ and $\beta = 2M/\sigma^2$, and
bisection on an explicit sign function finds $\xi$ (Proposition~\ref{prop:bisection}).
\item Three properties useful in practice: the optimal threshold scales with the carrying capacity, a constant
bias in the population estimate shifts it by the same factor, and no threshold strategy drives the population
extinct (Section~\ref{sec:remarks}).
\item Simulations of \eqref{eq:model} confirm the exact yield and illustrate the risk of low thresholds and the
loss from annual hunting seasons (Section~\ref{sec:numerics}).
\item For the logistic model, the uniqueness, existence and unimodality, and the statement of Conjecture~4.1
about the yield function, are proved in Lean~4 with Mathlib, with the yield written exactly as in
\cite[eq.~(4.4)]{hening19}; so is the abstract form of the uniqueness argument (Section~\ref{sec:lean}).
\end{itemize}

\section{Threshold strategies for one-dimensional diffusions}\label{sec:general}

Let the population follow
\begin{equation}\label{eq:general}
  \dd X_t = \bigl(f(X_t) - h_t\bigr)\dd t + v(X_t)\dd B_t ,\qquad 0 \le h_t \le k(X_t),
\end{equation}
on $(0,\infty)$, with continuous $f$, continuous $v > 0$ and continuous $k > 0$. Here $f$ is the total growth rate
of the population, and $k(x)$ the largest harvest per unit time at population size $x$; the logistic model has
$f(x) = x(\mu - \kappa x)$, $v(x) = \sigma x$ and $k(x) = Mx$. The threshold strategy with threshold $\eta > 0$
harvests $h_t = k(X_t)\mathbf 1\{X_t > \eta\}$. Fix $x_1 > 0$ and write
\[
  m(x) = \frac{1}{v(x)^2}\exp\Bigl(\int_{x_1}^x \frac{2f}{v^2}\Bigr),\qquad
  m_0(x) = m(x)\exp\Bigl(-\int_{x_1}^x \frac{2k}{v^2}\Bigr),\qquad
  \lambda(\eta) = \frac{m(\eta)}{m_0(\eta)} .
\]
Up to a constant factor, $m$ is the speed density of the unharvested process and $m_0$ that of the process
harvested at full capacity everywhere. The speed density of the threshold process is $m$ on $(0,\eta]$ and
$\lambda(\eta)m_0$ on $(\eta,\infty)$, since both are solutions of the same first-order equation and agree at
$\eta$. Let
\[
  A = \int_0^\eta m,\quad B = \int_\eta^\infty m_0,\quad P = \int_0^\eta f m,\quad
  Q = \int_\eta^\infty f m_0,\quad R = \int_\eta^\infty k m_0 ,
\]
all as functions of $\eta$.

\begin{assumption}\label{ass:general}
For every $\eta > 0$ the integrals $\int_0^\eta m$, $\int_0^\eta |f| m$, $\int_\eta^\infty m_0$,
$\int_\eta^\infty |f| m_0$ and $\int_\eta^\infty k m_0$ are finite, and $v^2 m \to 0$ as $x \to 0$ and
$v^2 m_0 \to 0$ as $x \to \infty$. There is $x_p \in [0,\infty]$ such that $f$ is strictly increasing on
$(0,x_p] \cap (0,\infty)$ and strictly decreasing on $[x_p,\infty) \cap (0,\infty)$.
\end{assumption}

Under Assumption~\ref{ass:general} the function
\begin{equation}\label{eq:H}
  H(\eta) = \frac{\lambda(\eta) R(\eta)}{A(\eta) + \lambda(\eta)B(\eta)} .
\end{equation}
is well defined. If the threshold process is positive recurrent, for instance when its scale function diverges at
$0$ and at $\infty$, its stationary density is proportional to $m$ on $(0,\eta]$ and to $\lambda(\eta)m_0$ on
$(\eta,\infty)$, and $H(\eta)$ is its stationary mean harvest and almost surely its long-run average yield, by the
ergodic theorem for one-dimensional diffusions; for the logistic model see \cite[Remark~2.2 and
eq.~(4.4)]{hening19}. The results below are statements about the function $H$.

\begin{lemma}[the yield is the mean growth]\label{lem:mean}
$H(\eta) = \dfrac{P + \lambda Q}{A + \lambda B}$.
\end{lemma}

\begin{proof}
By definition $(v^2 m)' = 2fm$ and $(v^2 m_0)' = 2(f - k)m_0$. Integrating over $(0,\eta)$ and $(\eta,\infty)$
and using the boundary conditions of Assumption~\ref{ass:general} gives $2P = v(\eta)^2 m(\eta)$ and
$2(Q - R) = -v(\eta)^2 m_0(\eta)$. Multiply the second identity by $\lambda$ and add it to the first; since
$\lambda m_0(\eta) = m(\eta)$, this gives $P + \lambda Q = \lambda R$.
\end{proof}

So the long-run yield equals the stationary mean of the total growth. Write
\[
  G = P/A, \qquad K = Q/B, \qquad D = K - G
\]
for the mean growth below and above the threshold, and their difference.

\begin{lemma}[sign of the derivative]\label{lem:sign}
$H$ is differentiable, and
$H'(\eta) = \lambda'(\eta)\,A B\,D(\eta)/(A + \lambda B)^2$ with $\lambda' = 2\lambda k/v^2 > 0$. Hence
$\sign H' = \sign D$.
\end{lemma}

\begin{proof}
$A' = m(\eta)$, $B' = -m_0(\eta)$, $P' = f(\eta)m(\eta)$ and $Q' = -f(\eta)m_0(\eta)$. Since $\lambda m_0(\eta) =
m(\eta)$, the terms from the moving endpoint cancel: $(P + \lambda Q)' = \lambda' Q$ and $(A + \lambda B)' =
\lambda' B$. The quotient rule gives the formula.
\end{proof}

\begin{lemma}[strict down-crossings]\label{lem:cross}
At every zero of $D$, $D' < 0$.
\end{lemma}

\begin{proof}
From the derivatives above, $G' = (m/A)(f - G)$ and $K' = (m_0/B)(K - f)$. At a zero of $D$ with common value
$G = K = c$,
\[
  D' = \bigl(c - f(\eta)\bigr)\Bigl(\frac{m_0(\eta)}{B} + \frac{m(\eta)}{A}\Bigr).
\]
If $\eta \le x_p$, then $f$ is strictly increasing on $(0,\eta]$, so its mean $G$ under $m$ is below $f(\eta)$.
If $\eta > x_p$, then $f$ is strictly decreasing on $[\eta,\infty)$, so its mean $K$ under $m_0$ is below
$f(\eta)$. In both cases $c < f(\eta)$ and $D' < 0$.
\end{proof}

A differentiable function whose derivative is negative at every zero has at most one zero: between two
consecutive zeros it would have to change sign from negative to positive. This down-crossing step is the
uniqueness argument in the proof of \cite[Lemma~5.4]{ekstrom26}, there under the assumptions of that paper.

\begin{theorem}\label{thm:general}
Under Assumption~\ref{ass:general}, the yield $H$ of \eqref{eq:H} has at most one critical point in $(0,\infty)$,
and a critical point is a strict local maximum. If $H$ attains its supremum over $(0,\infty)$ at $\eta^*$, then
$\eta^*$ is the only maximizer and the only critical point, $H$ is strictly increasing on $(0,\eta^*]$ and strictly
decreasing on $[\eta^*,\infty)$.
\end{theorem}

\begin{proof}
By Lemma~\ref{lem:sign} the critical points of $H$ are the zeros of $D$, and by Lemma~\ref{lem:cross} $D$ has at
most one zero, at which it changes sign from positive to negative. A maximizer is a critical point. If $\eta^*$ is
the zero of $D$, then $D$ has no zero on $(0,\eta^*)$ or on $(\eta^*,\infty)$, so by continuity $D > 0$ on the
first and $D < 0$ on the second, and the mean value theorem gives the monotonicity.
\end{proof}

\begin{remark}
If $f > 0$ on $(0,x_0)$ and $f < 0$ on $(x_0,\infty)$ for some $x_0$, then $P(x_0) > 0 > Q(x_0)$, so
$D(x_0) < 0$ and $H$ decreases at $x_0$. A maximizer then exists as soon as $H(\eta) < H(x_0)$ for some
$\eta < x_0$, by the intermediate value theorem applied to $D$.
\end{remark}

\section{The stochastic logistic model}\label{sec:logistic}

For \eqref{eq:model} we use the notation
\[
  s = \frac{2}{\sigma^2},\quad a = s\mu,\quad \beta = sM,\quad \alpha = a - \beta,\quad c = s\kappa,
\]
so that, up to constant factors, $m(y) = y^{a-2}e^{-cy}$, $m_0(y) = y^{\alpha-2}e^{-cy}$ and
$\lambda(\eta) = \eta^{\beta}$. The condition $\mu > \sigma^2/2$ is $a > 1$. Assumption~\ref{ass:general} holds:
$\int_0 m < \infty$ because $a > 1$, the tail integrals over $(\eta,\infty)$ are finite for every real $\alpha$
because of the exponential factor, $y^2 m = y^a e^{-cy} \to 0$ at $0$, $y^2 m_0 \to 0$ at $\infty$, and
$f(y) = y(\mu - \kappa y)$ increases up to $\mu/(2\kappa)$ and decreases after it. Multiplying both branches of
the stationary density \cite[eq.~(4.3)]{hening19} by the same constant shows that \eqref{eq:H} is exactly the
yield \cite[eq.~(4.4)]{hening19}. Note that $\alpha$ may be negative or zero: there is no condition on
$\mu - M$.

\begin{lemma}[a critical point exists]\label{lem:exist}
There is $\eta < \mu/\kappa$ with $D(\eta) > 0$, and $D(\mu/\kappa) < 0$.
\end{lemma}

\begin{proof}
Write $x_0 = \mu/\kappa$, the zero of $f$. By the remark above, $D(x_0) < 0$. For the first claim we use two
bounds, valid for every $\eta > 0$. First, $e^{-cy} \ge e^{-c\eta}$ on $(0,\eta)$ gives
$A \ge e^{-c\eta}\eta^{a-1}/(a-1)$, and $sP = \eta^a e^{-c\eta}$ from the proof of Lemma~\ref{lem:mean}; hence
$sG \le (a-1)\eta$. Second, $sQ = \beta R_1 - \eta^\alpha e^{-c\eta}$ with $R_1 = \int_\eta^\infty y\,m_0(y)\dd y$,
by the same proof.

\emph{Case $\alpha > 1$.} Then $B_0 = \int_0^\infty m_0 < \infty$, so $B(\eta) \le B_0$; also $R_1(\eta) \ge
R_1(1)$ for $\eta \le 1$, and $\eta^\alpha e^{-c\eta} \le \eta$. Hence
$sQ - (a-1)\eta B \ge \beta R_1(1) - \eta\bigl(1 + (a-1)B_0\bigr)$, which is positive for every
$\eta \le \min(1, x_0/2)$ with $\eta(1 + (a-1)B_0) \le \beta R_1(1)/2$. For such $\eta$,
$sK > (a-1)\eta \ge sG$.

\emph{Case $\alpha \le 1$.} By \eqref{eq:H}, $H(\eta) \le M R_1(\eta)/B(\eta)$. For $0 < \eta < \varepsilon$,
splitting the integral at $\varepsilon$ gives $R_1(\eta) \le \varepsilon B(\eta) + R_1(\varepsilon)$. For
$\eta < 1$, $y^{\alpha-2} \ge y^{-1}$ and $e^{-cy} \ge e^{-c}$ on $(\eta,1)$, so
$B(\eta) \ge e^{-c}\log(1/\eta)$, which tends to infinity. Take $\varepsilon \le H(x_0)/(2M)$; then for small
$\eta$ the bound gives $H(\eta) < H(x_0)$, and by the mean value theorem $H' > 0$, hence $D > 0$, somewhere in
$(\eta, x_0)$.
\end{proof}

\begin{theorem}[Conjecture~4.1 of \cite{hening19}]\label{thm:logistic}
For all $\mu > \sigma^2/2$, $\kappa > 0$ and $M > 0$, the yield $H$ of the threshold strategy in \eqref{eq:model}
has exactly one critical point $\eta^* \in (0,\infty)$. It lies below $\mu/\kappa$, $H$ is strictly increasing on
$(0,\eta^*]$ and strictly decreasing on $[\eta^*,\infty)$, and $\eta^*$ is the unique optimal threshold.
Harvesting at rate $M$ above $\eta^*$ and not below it is an optimal strategy among all admissible strategies.
\end{theorem}

\begin{proof}
Lemma~\ref{lem:exist} and the intermediate value theorem give a zero of $D$ in $(0,\mu/\kappa)$, and
Theorem~\ref{thm:general} gives the rest of the statement about $H$. By \cite[Theorem~2.1]{hening19} an
optimal strategy is a threshold strategy; its threshold maximizes $H$ and is therefore $\eta^*$.
\end{proof}

The case $\mu - M \le \sigma^2/2$ deserves a comment. There the population harvested at full rate everywhere
does not persist, $\int_0 m_0 = \infty$, and both mean growths $G$ and $K$ tend to $0$ as $\eta \to 0$, so a
comparison of their limits at $0$ decides nothing. The second case of Lemma~\ref{lem:exist} uses instead that
$H(\eta) \to 0$ as $\eta \to 0$ while $H(\mu/\kappa) > 0$.

\section{Closed form and an algorithm}\label{sec:closed}

Let $x = c\eta = 2\kappa\eta/\sigma^2$, and let $\gamma(\cdot,x)$ and $\Gamma(\cdot,x)$ be the lower and upper
incomplete gamma functions. Substituting $y = x'/c$ in $A$, $B$ and $R_1$ gives
\begin{equation}\label{eq:closed}
  H(\eta) = \frac{M\sigma^2}{2\kappa}\,h(x), \qquad
  h(x) = \frac{x^\beta\,\Gamma(\alpha,x)}{\gamma(a-1,x) + x^\beta\,\Gamma(\alpha-1,x)} .
\end{equation}
The function $\Gamma(s,x)$ is defined for every real $s$ when $x > 0$, so \eqref{eq:closed} covers $\alpha \le 1$.
It is elementary when $a$ and $\alpha$ are integers with $a, \alpha \ge 2$, and involves the exponential
integral when $\alpha$ is an integer below $2$. For a small cap ($\beta \to 0$ with $a > 1$ fixed), $h(\xi) - (\alpha - 1) \to 0$: the optimal yield
approaches $M(\mu - M - \sigma^2/2)/\kappa$, the yield of harvesting at the constant rate $M$.

With the recurrences $\gamma(s+1,x) = s\gamma(s,x) - x^se^{-x}$ and $\Gamma(s+1,x) = s\Gamma(s,x) + x^se^{-x}$, the
sign of $D$ becomes the sign of
\begin{equation}\label{eq:S}
  S(x) = \bigl(\beta\,\Gamma(\alpha,x) - x^\alpha e^{-x}\bigr)\gamma(a-1,x) - x^a e^{-x}\,\Gamma(\alpha-1,x).
\end{equation}
So $\eta^* = \sigma^2\xi/(2\kappa)$, where $\xi = \xi(a,\beta)$ is the unique zero of $S$. Equivalently
\begin{equation}\label{eq:ratio}
  \frac{\eta^*}{\mu/\kappa} = \frac{\xi(a,\beta)}{a} ,
\end{equation}
so the optimal threshold, as a fraction of the carrying capacity, depends only on $2\mu/\sigma^2$ and
$2M/\sigma^2$.

\begin{proposition}\label{prop:bisection}
$S(a) < 0$ and $S > 0$ on $(0,\xi)$. Let $j \ge 1$ be the first index with $S(a/2^j) > 0$, and start with
$\mathrm{lo} = a/2^j$ and $\mathrm{hi} = 2\,\mathrm{lo}$; then halve the bracket according to the sign of $S$ at
its midpoint. After $j + \lceil\log_2(1/\epsilon)\rceil$ evaluations of $S$, the bracket contains $\xi$ and
$(\mathrm{hi} - \mathrm{lo})/\mathrm{lo} \le \epsilon$.
\end{proposition}

\begin{proof}
$x = a$ corresponds to $\eta = \mu/\kappa$, where $D < 0$ by Lemma~\ref{lem:exist}; so $\xi < a$. By
Theorem~\ref{thm:logistic}, $D$ and hence $S$ are positive on $(0,\xi)$ and negative on $(\xi,\infty)$, so the
search for $j$ terminates, $S(\mathrm{hi}) \le 0$ by the choice of $j$, and bisection keeps $\xi$ inside the
bracket. Its initial relative width is $1$, the width halves at each step, and $\mathrm{lo}$ never decreases.
\end{proof}

Each evaluation of $S$ needs three incomplete gamma functions. For $\epsilon = 10^{-12}$ this means $40 + j$ evaluations, with $j \le 3$ in all cases of Table~\ref{tab:xi}.

\begin{table}[t]
\centering
\begin{tabular}{rrrrr}
\toprule
$a = 2\mu/\sigma^2$ & $\beta = 2M/\sigma^2$ & $\xi$ & $\eta^*/(\mu/\kappa)$ & $h(\xi)$\\
\midrule
1.25 & 1 & 0.492788 & 0.3942 & 0.082387\\
2 & 0.5 & 0.390607 & 0.1953 & 0.638548\\
2 & 2 & 0.928659 & 0.4643 & 0.281628\\
3 & 1 & 0.773581 & 0.2579 & 1.175154\\
5 & 2 & 1.564222 & 0.3128 & 2.207911\\
10 & 4 & 3.374603 & 0.3375 & 5.130319\\
\bottomrule
\end{tabular}
\caption{Optimal scaled threshold $\xi$, the threshold as a fraction of the carrying capacity, and the scaled
optimal yield $h(\xi) = H^*/(M\sigma^2/(2\kappa))$. The row $a = \beta = 2$ is the example of Figure~1 of
\cite{hening19}.}
\label{tab:xi}
\end{table}

\paragraph{Worked example.}
Measure the population in units of its carrying capacity, so that $\kappa = \mu$. A growth rate of
$\mu = 0.14$ per year was estimated for a reintroduced ibex population \cite{saether02}, and an environmental
variance of $\sigma^2 = 0.007$ per year for a harvested wild reindeer population \cite{kvalnes24}. With these
values, and with the larger variances $0.02$ and $0.05$ as illustrative cases, Table~\ref{tab:example} lists the
optimal threshold and yield for several caps. The deterministic optimum of the capped logistic model is
$M(1 - M/\mu)$ for $M \le \mu/2$ and $\mu/4$ for $M \ge \mu/2$, attained at the population level
$\max(1 - M/\mu,\, 1/2)$. With noise the optimal
threshold in these examples lies below half the carrying capacity and approaches one half as the cap grows.

\begin{table}[t]
\centering
\begin{tabular}{rrrrrr}
\toprule
$\sigma^2$ & $M$ & $\alpha$ & $\eta^*$ & $H^*$ & deterministic optimum\\
\midrule
0.007 & 0.05 & 25.71 & 0.3417 & 0.03090 & 0.03214\\
0.007 & 0.10 & 11.43 & 0.4787 & 0.03469 & 0.03500\\
0.007 & 0.14 & 0 & 0.4992 & 0.03486 & 0.03500\\
0.02 & 0.05 & 9 & 0.3212 & 0.02872 & 0.03214\\
0.02 & 0.10 & 4 & 0.4583 & 0.03350 & 0.03500\\
0.05 & 0.05 & 3.6 & 0.2937 & 0.02436 & 0.03214\\
0.05 & 0.10 & 1.6 & 0.4324 & 0.02995 & 0.03500\\
0.05 & 0.14 & 0 & 0.4877 & 0.03127 & 0.03500\\
\bottomrule
\end{tabular}
\caption{Optimal threshold and long-run yield per year for $\mu = \kappa = 0.14$ (population in units of the
carrying capacity). The value $\mu = 0.14$ is taken from \cite{saether02} and $\sigma^2 = 0.007$ from
\cite{kvalnes24}; the other variances are illustrative. The rows with $\alpha = 0$ have $M = \mu$, where
harvesting at full rate everywhere drives the population extinct.}
\label{tab:example}
\end{table}

\section{Three properties for practice}\label{sec:remarks}

\begin{proposition}[scaling]\label{prop:scale}
For fixed $a$ and $\beta$, the optimal threshold is proportional to the carrying capacity $\mu/\kappa$, by
\eqref{eq:ratio}. If the carrying capacity is misjudged by a factor $\theta$, the threshold used is
$\theta\eta^*$ and the yield is $H(\theta\eta^*)$; the relative loss $1 - h(\theta\xi)/h(\xi)$ depends only on
$\theta$, $a$ and $\beta$.
\end{proposition}

\begin{proposition}[biased population estimates]\label{prop:bias}
Suppose decisions use an estimate $\hat X_t = \theta X_t$ with a constant factor $\theta > 0$. A threshold $b$ on
$\hat X$ is the threshold $b/\theta$ on $X$, with yield $H(b/\theta)$. The best threshold on the estimate is
$\theta\eta^*$, and it achieves the optimal yield.
\end{proposition}

Both statements follow directly from the definitions and from \eqref{eq:closed}. Random errors in the population
estimate are a different matter; see Section~\ref{sec:scope}.

\begin{proposition}[persistence]\label{prop:persist}
For every threshold $\eta > 0$ the harvested population in \eqref{eq:model} is positive recurrent on
$(0,\infty)$; in particular it neither goes extinct nor explodes. This holds also when $\mu - M \le \sigma^2/2$.
\end{proposition}

\begin{proof}
The speed density of the threshold process is integrable on $(0,\infty)$, since $\int_0 m < \infty$ for $a > 1$
and the tail decays exponentially. Its scale density is proportional to $y^{-a}e^{cy}$ near $0$ and to
$y^{-\alpha}e^{cy}$ near $\infty$, so the scale function diverges at both ends, and both boundaries are
inaccessible \cite[Chapter~15]{karlin81}; in fact $0$ is natural and $\infty$ is entrance
\cite[Section~4]{hening19}. A regular diffusion with inaccessible boundaries and finite speed
measure is positive recurrent.
\end{proof}

Below the threshold the process is the unharvested one, which is pushed away from $0$ when $\mu > \sigma^2/2$.
In the model, a low threshold therefore does not cause extinction, but it does make low population levels
more frequent; Figure~\ref{fig:risk} shows the probability of falling below one percent of the carrying capacity
within ten time units.

\section{Simulations}\label{sec:numerics}

We simulate \eqref{eq:model} with the parameters of Figure~1 of \cite{hening19}, $\mu = M = \kappa = \sigma^2 = 1$,
for which $\eta^* = 0.4643297$ and $H(\eta^*) = 0.1408139$. To keep the population positive we use the Euler
scheme for $Y = \log X$,
\[
  Y_{t+\Delta} = Y_t + \bigl(\mu - \kappa e^{Y_t} - M\mathbf 1\{e^{Y_t} > \eta\} - \sigma^2/2\bigr)\Delta
  + \sigma\sqrt{\Delta}\,Z_t ,
\]
with independent standard normal $Z_t$ and step $\Delta = 0.01$. On the same Brownian paths, halving the step
changes the time-averaged yield at $\eta^*$ by $0.14\%$ (averages over 180 time units after a burn-in of 20,
100 paths). All runs use fixed seeds,
and the scripts check their results with assertions.

Figure~\ref{fig:paths} shows sample paths at $\eta^*/2$, $\eta^*$ and $2\eta^*$ on the same noise. With $\sigma^2 = 1$ the noise is strong, and all
three paths spend long periods far below the carrying capacity. The low threshold harvests in every recovery
and keeps the population small, while the high threshold harvests only at the peaks. Figure~\ref{fig:yield}
compares the exact yield \eqref{eq:closed} with time averages over $270$ time units after a burn-in of $30$,
for 200 paths per threshold. All seven estimates agree with the exact value within four standard errors plus
$0.003$ for the bias of the time step.

\begin{figure}[t]
\centering
\begin{tikzpicture}
\begin{groupplot}[group style={group size=1 by 3, vertical sep=0.55cm, xticklabels at=edge bottom},
  width=0.92\linewidth, height=3.3cm, xmin=0, xmax=30, ymin=0, ymax=2.2, ylabel={$X_t$}]
\nextgroupplot[title={\footnotesize threshold $\eta^*/2 = 0.232$}, title style={yshift=-1.5ex}]
\addplot[harvest, const plot] table[col sep=comma, x=t, y expr=2.2*\thisrow{harvest}] {\datadir/paths_low.csv} \closedcycle;
\addplot[pathline] table[col sep=comma, x=t, y=x] {\datadir/paths_low.csv};
\addplot[thr, domain=0:30] {0.2322};
\nextgroupplot[title={\footnotesize optimal threshold $\eta^* = 0.464$}, title style={yshift=-1.5ex}]
\addplot[harvest, const plot] table[col sep=comma, x=t, y expr=2.2*\thisrow{harvest}] {\datadir/paths_opt.csv} \closedcycle;
\addplot[pathline] table[col sep=comma, x=t, y=x] {\datadir/paths_opt.csv};
\addplot[thr, domain=0:30] {0.4643};
\nextgroupplot[title={\footnotesize threshold $2\eta^* = 0.929$}, title style={yshift=-1.5ex}, xlabel={time $t$}]
\addplot[harvest, const plot] table[col sep=comma, x=t, y expr=2.2*\thisrow{harvest}] {\datadir/paths_high.csv} \closedcycle;
\addplot[pathline] table[col sep=comma, x=t, y=x] {\datadir/paths_high.csv};
\addplot[thr, domain=0:30] {0.9287};
\end{groupplot}
\end{tikzpicture}

\caption{Sample paths of \eqref{eq:model} with $\mu = M = \kappa = \sigma^2 = 1$ under three thresholds (dashed),
driven by the same noise. Shaded: periods of harvesting.}
\label{fig:paths}
\end{figure}

\begin{figure}[t]
\centering
\begin{tikzpicture}
\begin{axis}[xmin=0, xmax=2.5, ymin=0, ymax=0.16, xlabel={threshold $\eta$}, ylabel={long-run yield},
  height=5cm, legend pos=north east, legend cell align=left, scaled y ticks=false, yticklabel style={/pgf/number format/fixed}]
\addplot[black, thick] table[col sep=comma, x=b, y=H] {\datadir/yield_exact.csv};
\addlegendentry{exact $H(\eta)$}
\addplot[only marks, mark=o, blue!70!black, error bars/.cd, y dir=both, y explicit]
  table[col sep=comma, x=b, y=H, y error=ci] {\datadir/yield_mc.csv};
\addlegendentry{simulation, 95\% interval}
\draw[thr] (axis cs:0.4643,0) -- (axis cs:0.4643,0.16);
\node[font=\footnotesize, anchor=south west, red!70!black] at (axis cs:0.48,0.005) {$\eta^*$};
\end{axis}
\end{tikzpicture}

\caption{Exact long-run yield $H(\eta)$ and simulated time averages with 95\% confidence intervals, for
$\mu = M = \kappa = \sigma^2 = 1$.}
\label{fig:yield}
\end{figure}

Figure~\ref{fig:sens} shows the optimal threshold and yield for $\mu = \kappa = 1$ as functions of the cap $M$ for
five noise levels, computed with Proposition~\ref{prop:bisection}. The dot on each curve marks $M = \mu -
\sigma^2/2$; to its right lies the regime that is outside the assumptions of \cite[Lemma~5.4]{ekstrom26}. The
curves pass through it smoothly. On this grid, larger noise lowers the optimal yield at every cap. Its effect on the
optimal threshold depends on the cap: it lowers $\eta^*$ for $M \le 1.2$, raises it for $M \ge 1.9$, and is not
monotone in between.

\begin{figure}[t]
\centering
\begin{tikzpicture}
\begin{groupplot}[group style={group size=2 by 1, horizontal sep=1.6cm}, width=0.5\linewidth, height=5cm,
  xmin=0, xmax=2.05, xlabel={harvesting cap $M$}, cycle list={{blue!80!black},{teal},{black},{orange!80!black},{red!70!black}}]
\nextgroupplot[ylabel={optimal threshold $\eta^*$}, ymin=0, ymax=0.65, legend to name=senslegend, legend columns=5]
\foreach \s in {0.2,0.6,1.0,1.4,1.8} {
  \addplot+[thick, mark=none] table[col sep=comma, x=M, y=eta] {\datadir/sens_\s.csv};
  \addlegendentryexpanded{$\sigma^2 = \s$}
}
\foreach \s/\m in {0.2/0.9,0.6/0.7,1.0/0.5,1.4/0.3,1.8/0.1} {
  \addplot[only marks, mark=*, mark size=1.4pt] table[col sep=comma, x=M, y=eta, restrict expr to domain={abs(\thisrow{M}-\m)}{0:0.001}] {\datadir/sens_\s.csv};
}
\nextgroupplot[ylabel={optimal yield $H^*$}, ymin=0, ymax=0.26, scaled y ticks=false, yticklabel style={/pgf/number format/fixed}]
\foreach \s in {0.2,0.6,1.0,1.4,1.8} {
  \addplot+[thick, mark=none] table[col sep=comma, x=M, y=H] {\datadir/sens_\s.csv};
}
\end{groupplot}
\node at ($(group c1r1.south)!0.5!(group c2r1.south)+(0,-1.2cm)$) {\pgfplotslegendfromname{senslegend}};
\end{tikzpicture}

\caption{Optimal threshold $\eta^*$ (left) and optimal yield $H^*$ (right) against the cap $M$, for $\mu = \kappa = 1$
and five values of $\sigma^2$. Dots: $M = \mu - \sigma^2/2$.}
\label{fig:sens}
\end{figure}

\begin{figure}[t]
\centering
\begin{tikzpicture}
\begin{axis}[xmin=0, xmax=2.1, ymin=0, ymax=0.5, xlabel={threshold $\eta$},
  ylabel={$\mathbb{P}(\min_{t\le 10} X_t < 0.01)$}, height=4.6cm]
\addplot[blue!70!black, mark=*, mark size=1.2pt, error bars/.cd, y dir=both, y explicit]
  table[col sep=comma, x=b, y=p, y error=ci] {\datadir/risk.csv};
\draw[thr] (axis cs:0.4643,0) -- (axis cs:0.4643,0.5);
\node[font=\footnotesize, anchor=south west, red!70!black] at (axis cs:0.48,0.4) {$\eta^*$};
\end{axis}
\end{tikzpicture}

\caption{Probability that the population falls below $0.01$ within ten time units, starting at the carrying
capacity $1$, against the threshold; $2000$ paths per threshold, 95\% intervals, $\mu = M = \kappa = \sigma^2 = 1$.
Monitoring at the time step slightly underestimates the probability.}
\label{fig:risk}
\end{figure}

\paragraph{Annual hunting seasons.}
Hunting usually takes place in seasons, not continuously. In a simple seasonal variant the population grows
without harvest during each year, and at the end of the year, if $X > \eta$, either the fraction
$u = 1 - e^{-M}$ of the population is taken (the fraction a constant rate $M$ would remove in one year without
growth), or the surplus $X - \eta$, at most $uX$. Figure~\ref{fig:seasons} compares the resulting long-run yields
per year (200 paths, averages over 270 years after a burn-in of 30) with the continuous optimum. At their best
thresholds on the grid the seasonal rules lose about 13.5 percent (proportional rule) and 9.5 percent (surplus rule)
of the continuous optimal yield. The best threshold of the
proportional rule lies near $0.5$, slightly above $\eta^*$, and that of the surplus rule near $0.3$. The theory of Section~\ref{sec:general} does not cover these discrete rules.

\begin{figure}[t]
\centering
\begin{tikzpicture}
\begin{axis}[xmin=0, xmax=2.1, ymin=0, ymax=0.16, xlabel={threshold $\eta$}, ylabel={long-run yield per year},
  height=5cm, legend pos=north east, legend cell align=left, scaled y ticks=false, yticklabel style={/pgf/number format/fixed}]
\addplot[black, thick] table[col sep=comma, x=b, y=H, restrict x to domain=0:2.1] {\datadir/yield_exact.csv};
\addlegendentry{continuous harvesting, exact}
\addplot[blue!70!black, mark=o, mark size=1.3pt, error bars/.cd, y dir=both, y explicit]
  table[col sep=comma, x=b, y=prop, y error=prop_ci] {\datadir/seasons.csv};
\addlegendentry{annual season, fraction $u$}
\addplot[orange!80!black, mark=triangle, mark size=1.5pt, error bars/.cd, y dir=both, y explicit]
  table[col sep=comma, x=b, y=surplus, y error=surplus_ci] {\datadir/seasons.csv};
\addlegendentryexpanded{annual season, surplus, at most $u X$}
\end{axis}
\end{tikzpicture}

\caption{Long-run yield per year of two annual-season rules against the threshold, with 95\% intervals, and the
exact yield of continuous harvesting, for $\mu = M = \kappa = \sigma^2 = 1$.}
\label{fig:seasons}
\end{figure}

\section{Scope}\label{sec:scope}

The results concern one population in continuous time with a known model. Several features of real game
management lie outside this setting.
\begin{itemize}
\item \emph{Seasons and annual quotas.} Discrete-time harvesting has no explicit stationary density, and the
proofs above do not apply; for discrete-time stochastic population models see \cite{hening21}.
\item \emph{Age and sex structure.} Cull plans usually prescribe numbers per age and sex class. A structured
population is a multidimensional diffusion, again without an explicit stationary density.
\item \emph{Movement between hunting districts.} Two coupled populations with separate owners lead to a
two-dimensional problem and, with separate objectives, to a game.
\item \emph{Parameter uncertainty and random measurement error.} Proposition~\ref{prop:scale} describes the
effect of a misjudged carrying capacity. An expected yield over a distribution of parameters is a mixture of
unimodal functions and need not be unimodal. For random errors in the population estimate, the optimal strategy
need not be of threshold type; numerical studies are in \cite{memarzadeh18}.
\end{itemize}

\section{Formal verification}\label{sec:lean}

\begin{sloppypar}
The following statements are proved in Lean~4 (toolchain v4.36.0-rc1) with Mathlib (revision
\texttt{d875a8de}), in the folder \texttt{laif/bounded-harvesting-threshold/lean} of
\url{https://github.com/ssenge/LAIF}. Every theorem depends only on the standard axioms \texttt{propext},
\texttt{Classical.choice} and \texttt{Quot.sound}.
\end{sloppypar}
\begin{itemize}
\item \texttt{Crossing.lean}: a function on $(0,\infty)$ whose derivative is negative at every zero has at most
one zero; the derivative formula of Lemma~\ref{lem:sign}; and the abstract form of Theorem~\ref{thm:general}
(\texttt{critical\_point\_unique}): for differentiable $A, B, P, Q, \lambda$ with the derivative relations of
Lemma~\ref{lem:sign}, positive $A, B, \lambda, \lambda', m, m_0$, $m = \lambda m_0$, and $G < f$ or $K < f$ at every
point, the quotient $(P + \lambda Q)/(A + \lambda B)$ has at most one critical point.
\item \texttt{Model.lean}: the logistic model with the yield defined literally as \cite[eq.~(4.4)]{hening19},
including the normalizing constant of the density \cite[eq.~(4.3)]{hening19}; Lemma~\ref{lem:mean}; the
derivative of the yield; uniqueness of the critical point for all parameters
(\texttt{yieldH\_critical\_point\_unique}).
\item \texttt{General.lean}: the general model of Section~\ref{sec:general} with $m$, $m_0$ and $\lambda$
defined from $f$, $v$, $k$, Assumption~\ref{ass:general}, Lemma~\ref{lem:mean}, the derivative of $H$ and the
uniqueness part of Theorem~\ref{thm:general} (\texttt{Diffusion.Assumption.critical\_point\_unique}); an example
with drift $1 - x$, noise $\sqrt{2x}$ and cap $x/2$ satisfies all of Assumption~\ref{ass:general}, so its
hypotheses can hold together (\texttt{Diffusion.Example.nonvacuous}).
\item \texttt{Existence.lean}: Lemma~\ref{lem:exist}, strict monotonicity on both sides of the critical point,
and Conjecture~4.1 of \cite{hening19} in the form: there is exactly one $x^* > 0$ with $H'(x^*) = 0$, and it
maximizes $H$ over all thresholds (\texttt{conjecture\_4\_1\_unconditional}), for all $\mu > \sigma^2/2$,
$\kappa > 0$, $M > 0$, with no further hypothesis.
\end{itemize}
The statement that the threshold strategy is optimal among all admissible strategies uses
\cite[Theorem~2.1]{hening19} and is not formalized, nor is the identity between \eqref{eq:H} and the long-run
average yield. The second part of Theorem~\ref{thm:general}, the closed form \eqref{eq:closed},
Proposition~\ref{prop:bisection} and Section~\ref{sec:remarks} are proved on paper. The closed form, the
bisection and the simulations are checked numerically by the scripts in \texttt{sim/} of the same folder, which
also confirm unimodality of $H$ on 400 random parameter sets, 188 of them with $\mu - M \le \sigma^2/2$.

\section{Use of AI}
Claude Opus~5.5 (Anthropic) was used throughout this work: to explore approaches, to develop and check the
proofs, to write the Lean formalization and the scripts, and to draft the text.

\begin{thebibliography}{99}
\bibitem{hening19} A. Hening, D. H. Nguyen, S. C. Ungureanu, T. K. Wong. Asymptotic harvesting of populations in
random environments. \emph{J. Math. Biol.} 78 (2019) 293--329. arXiv:1710.01221.
\bibitem{ekstrom26} E. Ekstr\"om, D. Mata, H. Saarinen. Resource extraction from a stochastic population with
harvesting quotas. arXiv:2608.17778 (2026).
\bibitem{lande95} R. Lande, S. Engen, B.-E. S\ae ther. Optimal harvesting of fluctuating populations with a risk
of extinction. \emph{Am. Nat.} 145(5) (1995) 728--745.
\bibitem{lande97} R. Lande, B.-E. S\ae ther, S. Engen. Threshold harvesting for sustainability of fluctuating
resources. \emph{Ecology} 78(5) (1997) 1341--1350.
\bibitem{alvarez98} L. H. R. Alvarez, L. A. Shepp. Optimal harvesting of stochastically fluctuating populations.
\emph{J. Math. Biol.} 37(2) (1998) 155--177.
\bibitem{alvarezhening22} L. H. R. Alvarez E., A. Hening. Optimal sustainable harvesting of populations in random
environments. \emph{Stoch. Proc. Appl.} 150 (2022) 678--698.
\bibitem{heningtran20} A. Hening, K. Tran. Harvesting and seeding of stochastic populations: analysis and
numerical approximation. \emph{J. Math. Biol.} 81 (2020) 65--112. arXiv:1911.00734.
\bibitem{hening21} A. Hening. Coexistence, extinction, and optimal harvesting in discrete-time stochastic
population models. \emph{J. Nonlinear Sci.} 31 (2021). arXiv:1911.08398.
\bibitem{memarzadeh18} M. Memarzadeh, C. Boettiger. Resolving the measurement uncertainty paradox in ecological
management. arXiv:1812.11184 (2018).
\bibitem{saether02} B.-E. S\ae ther, S. Engen, F. Filli, R. Aanes, W. Schr\"oder, R. Andersen. Stochastic
population dynamics of an introduced Swiss population of the ibex. \emph{Ecology} 83(12) (2002) 3457--3465.
\bibitem{kvalnes24} T. Kvalnes, \O. Flagstad, J. V\aa ge, O. Strand, H. Viljugrein, B.-E. S\ae ther. Harvest and
decimation affect genetic drift and the effective population size in wild reindeer. \emph{Evol. Appl.} 17(4)
(2024) e13684.
\bibitem{karlin81} S. Karlin, H. M. Taylor. \emph{A Second Course in Stochastic Processes}. Academic Press, 1981.
\bibitem{bjagdg} Bundesjagdgesetz, \S~21 Abs.~2. \url{https://www.gesetze-im-internet.de/bjagdg/__21.html}.
\bibitem{djv} Deutscher Jagdverband. Jahresstrecken, hunting year 2024/25. Handbuch 2026.
\end{thebibliography}

\end{document}
